% Part: sets-functions-relations
% Chapter: induction
% Section: induction

\documentclass[../../include/open-logic-section]{subfiles}

\begin{document}

\olfileid{sfr}{ind}{int}

\olsection{引言}

\begin{explain}
归纳是逻辑学中常用的数学证明方法。你可能已经在数学中见过它，在那里你在自然数上做归纳。
事实证明，有许多不同种类的数学与逻辑对象集合都适合进行归纳。

一般而言，归纳法允许我们证明一个全称论断；也就是说，证明某一类对象中的每一个都具有某种性质。
特别地，当“全称引入”的证明方法过于复杂时，归纳法通常非常有用。
归纳法只适用于以特殊方式构造的数学对象：如果一个集合中的每个元素或者是基本元素，或者是由基本元素构造而来的，那么我们就可以对它使用归纳法。更准确地说：
\end{explain}

\begin{defn}
集合 $S$ \emph{对}函数 $f$ \emph{封闭}，当且仅当若 $x \in S$，则 $f(x) \in S$。
\end{defn}

\begin{explain}
例如，自然数集（$\mathbb{N}$）对后继函数 $f(x) = x+1$ 封闭。如果 $x$ 是自然数，那么 $x+1$ 也是自然数。然而，自然数集对前驱函数 $f(x) = x-1$ 不封闭，因为 $0$ 是自然数，但 $-1$ 不是。
\end{explain}

\begin{defn}
称性质 $P$ \emph{在}函数 $f(x)$ \emph{下保持}，是指对 $f$ 定义域中的任意对象 $a$，$P(a)$ 蕴含 $P(f(a))$（即若 $a$ 具有性质 $P$，则 $f(a)$ 也具有性质 $P$）。
\end{defn}

\begin{explain}
“是偶数”这一性质在函数 $f(x) = x+2$ 下保持，因为若 $x$ 是偶数，则 $x+2$ 也是偶数。该性质在后继函数下不保持，因为若 $x$ 是偶数，则 $x+1$ 是奇数。

归纳法对自然数有效，是因为每个自然数都可以通过从 $0$ 开始有限次应用后继函数而得到。任何我们可以对其做归纳的集合都具有这一特征。
\end{explain}

\begin{defn}[自然数上的归纳]
如果 $0$ 具有某个性质 $P$，且 $P$ 在后继函数下保持，那么每个自然数都具有性质 $P$。
\end{defn}

\begin{ex} 前 $n$ 个自然数的和为 $\frac{n(n+1)}{2}$。 \\

\textbf{基础步骤.} 前 $0$ 个自然数的和为 $0=
\frac{0\cdot 1}{2}$。\\

\textbf{归纳步骤.} 假设该性质对 $k$ 成立。我们将证明它对 $k+1$ 也成立。换句话说，假设 $P(k)$ 成立，即
\[ 0 + 1 + \ldots + k = \frac{k(k+1)}{2} \]
利用 $P(k)$ 的假设，我们将证明 $P(k+1)$ 成立，即
\[ 0 + 1 + \ldots + k + (k+1) = \frac{(k+1)(k+2)}{2} \]

\begin{align*}
0 + 1 + \ldots + k &= \frac{k(k+1)}{2} \tag{Assumption} \\
0 + 1 + \ldots + k + (k+1) &= \frac{k(k+1)}{2} + (k+1)
\tag{Add $k+1$ to both sides} \\
&= \frac{k(k+1) + 2(k+1)}{2} \\
&= \frac{2k + k + 2k +2}{2} \\
&=\frac{(k+1)(k+2)}{2}
\end{align*}
因此，归纳步骤成立，从而证明了该论断。
\end{ex}

\begin{explain}
对于公式，基本元素是原子公式。更复杂的公式通过使用（逻辑）算子连接较简单的公式而得到。每个算子都可以看作对应一个以公式为自变元的函数，该函数返回一个用该算子将它们连接而成的新公式。例如，将两个公式 $!!^a$ 和 $!B$ 连接成合取式的函数可以写作 $\mathcal E _\land (!!^a,
!B) = !!^a \land !B$。我们可以称这些为\emph{公式构造函数}。公式的集合对这些函数封闭：若 $!!^a$ 和 $!B$ 是公式，则 $\lnot !!^a$、$!!^a \land !B$ 等也是公式。
\end{explain}

\begin{defn}[公式上的归纳]
如果每个原子公式都具有性质 $P$，且 $P$ 在公式构造函数下保持，那么每个公式都具有性质 $P$。
\end{defn}

\begin{explain}
这个定义为公式上的归纳证明提供了一个“配方”。若要证明每个公式都具有性质 $P$，请遵循以下步骤：\\

\textbf{基础步骤.} 设 $!!^a$ 是一个原子公式。[\ldots] 因此，$!!^a$ 具有性质 $P$。

\textbf{归纳步骤.} 设 $!!^a$ 和 $!B$ 都是公式，且二者都具有性质 $P$。

$\lnot$ 情形：[\ldots] 因此，$\lnot !!^a$ 具有性质 $P$。

$\land$ 情形：[\ldots] 因此，$!!^a \land !B$ 具有性质 $P$。

$\lor$ 情形：[\ldots] 因此，$!!^a \lor !B$ 具有性质 $P$。

$\lif$ 情形：[\ldots] 因此，$!!^a \lif !B$ 具有性质 $P$。

$\liff$ 情形：[\ldots] 因此，$!!^a \liff !B$ 具有性质 $P$。

$\lforall[]$ 情形：[\ldots] 因此，$\lforall[x] !!^a$ 具有性质 $P$。

$\lexists[]$ 情形：[\ldots] 因此，$\lexists[x] !!^a$ 具有性质 $P$。 \\

因此，每个公式都具有性质 $P$。
\end{explain}

\end{document}